\section{Discussion}\label{sec:discussion}

\subsection{Lean explains the ladder}

Lean thinking faced a version of the problem this paper faces: work done quickly by many hands,
more of it than anyone could inspect, where a defect that travelled downstream cost far more to
fix than one caught where it was made \citep{ohno1988tps,womack1996lean}. Its answer was to build
the specification into the work itself. We read the Assertion Ladder as that answer applied to
a repository and a memoryless contributor, and this section makes the case through four Lean
ideas: value, standard work, the value stream and root cause analysis. Table~\ref{tab:Lean} maps each to the
receivers and to the evidence, and the section closes with the places where the analogy stops.

\begin{table*}[t]
  \centering\small
  \caption{Lean ideas, the rungs and patterns that carry them, and where they appear in the
  receivers. The first four rows are the argument; the last three are supporting mechanisms.}
  \label{tab:Lean}
  \begin{tabularx}{\linewidth}{>{\raggedright\arraybackslash}p{2.5cm}|>{\raggedright\arraybackslash}p{2.9cm}|X|X}
    \hline
    \textbf{Lean idea} & \textbf{Rung or pattern} & \textbf{In the receivers} & \textbf{Evidence} \\
    \hline\hline
    Value, and waste removed & fading; \pattern{Short Manifest} & no DI container, no ORM, one entity; no prose for what a test checks & 637 / 729 / 957 $\to$ 283 / 223 / 242 lines; manifests 480 / 427 / 468 words \\
    \hline
    Standard work & L1 \pattern{Worked Slice}, L2 \pattern{Short Manifest}, with L3 as its built-in test & the slice and the feature recipe, followed by every contributor and changed only by people & every pilot run followed the slice's layering \\
    \hline
    Value stream: defects found earlier & L3--L5 & V8 and V10 now stopped at the gate, not left to review; V7 moved from a behavioural test to an architecture rule in Python and TypeScript & V8 and V10 passed the gate in all three receivers at \texttt{8c4a74a}; caught now \\
    \hline
    Root cause analysis and kaizen & \pattern{Executable Rule}, \pattern{Teaching Failure}; \pattern{Honest Double} & the misses traced to one cause and closed by two rules; the V9 fixture corrected & 30 / 36 $\to$ 42 / 42; V9 caught in all three receivers \\
    \hline
    Jidoka (the line stops itself) & \pattern{Single Gate} & one verify command; ``change the code, not the test'' & no pilot run modified a rule (18 / 18) \\
    \hline
    Andon (signal and cord) & L4 \pattern{Teaching Failure}; \pattern{Human Boundary} & failures name the rule and the fix; ``ask a human first'' for dependencies, rules and schema & message on every rule; boundary declared (L2), not yet enforced \\
    \hline
    Poka-yoke (mistake-proofing) & L5 \pattern{Unrepresentable Mistake} & exhaustive error mapping that will not compile when a case is missing & V6 caught by the compiler in Java and TypeScript \\
    \hline
  \end{tabularx}
\end{table*}

\paragraph{Value} Lean asks first what the customer values, and the ladder adds the domain to the
answer: a change is valuable only if it works and leaves the domain's rules in one place. Sprawl is
then easy to name as waste. It is work that looks finished but will be redone, and the rework
falls on whoever next tries to change the rule. The same lens explains the force most easily
forgotten, ease of use. A paragraph of prose that repeats what a test enforces is read by every
contributor in every session and adds nothing the gate does not already guarantee; for a
contributor that relearns everything each time, that reading is the dominant cost. Fading removes
it. The receivers show both halves: application code fell from 637, 729 and 957 lines to 283, 223
and 242 when containers, mappers and a second entity were removed, and the manifests stayed under
500 words while two new kinds of intent were added.

\paragraph{Standard work} Standard work is where the mapping is tightest. The worked slice and the
feature recipe of Section~\ref{sec:receivers} (state the rule in EARS, let its pattern choose the
layer, copy the slice, run the gate) specify the content of a change, its sequence and its
expected outcome, which is how Spear and Bowen describe Toyota's work \citep{spear1999decoding}.
The ladder then does what they found Toyota does: it attaches a test to the standard, so that a
departure shows up at once. The lower rungs show and tell the standard; the upper rungs make it
self-checking. Every pilot run followed the slice's layering, and the two departures the pilot
recorded both came from the condition with the slice alone and no written standard. Standard work
is also owned. In Lean, the people who do the work improve the standard; in the receivers, an agent
follows the standard, and the \pattern{Human Boundary} requires a person's decision before the
rules or the manifest change. That division is deliberate, and we return to it below.

\paragraph{The value stream} Figure~\ref{fig:valuestream} draws the value stream of one change:
written by an agent, verified by the single gate, reviewed by a person, merged and released. Each
rung decides where along that stream a broken rule is found. At L0 only a reviewer who reads
closely can see it, and cognitive erosion is the loss of that reading. At L1 and L2 the
standard makes the mistake less likely but finds nothing. At L3 and L4 the gate finds it before
any person has spent time on the change, and at L5 it cannot be written. The seeding results show
the shift. At commit \texttt{8c4a74a}, V8 and V10 passed the gate in all three receivers and would
have reached review; now both are stopped at the gate by an architecture rule. V7 moved earlier
within the gate, from a behavioural test to an architecture rule, in Python and TypeScript, and V6
was already refused by the compiler in Java and TypeScript. We have no timing data for the stream, so we claim
only the position at which a defect is found, not the time saved. The Lean argument is that
position matters: a defect found at the gate goes back to the contributor who made it while the
change is still open, whereas one found in review has already cost a person's attention and, if
missed, becomes the rework that sprawl describes.

\begin{figure*}[t]
  \centering
  \resizebox{0.92\linewidth}{!}{% The value stream of one change, and where each rung finds a defect along it.
\begin{tikzpicture}[
    font=\footnotesize,
    step/.style={draw=#1, line width=0.9pt, fill=#1!7, rounded corners=2pt, text width=2.2cm,
                 minimum height=1.1cm, align=center, inner sep=4pt},
    arr/.style={-{Stealth[length=2.3mm]}, line width=0.8pt, #1},
    note/.style={text width=2.3cm, align=center, font=\scriptsize, anchor=north},
  ]
  % the stream
  \node[step=structure] (write) {\textbf{Write}\\an agent edits};
  \node[step=behaviour, right=0.55cm of write] (verify) {\textbf{Verify}\\the single gate};
  \node[step=missed, right=0.55cm of verify] (review) {\textbf{Review}\\a person reads};
  \node[step=muted, right=0.55cm of review] (merge) {\textbf{Merge}};
  \node[step=muted, right=0.55cm of merge] (prod) {\textbf{Release}};
  \draw[arr=muted] (write) -- (verify);
  \draw[arr=muted] (verify) -- (review);
  \draw[arr=muted] (review) -- (merge);
  \draw[arr=muted] (merge) -- (prod);
  % where each rung acts
  \node[note] at ($(write.south)+(0,-0.15)$) {L1, L2 guide the change\\L5 refuses to compile it\\[2pt]\textcolor{muted}{V6 in Java and TypeScript}};
  \node[note] at ($(verify.south)+(0,-0.15)$) {L3 finds the defect\\L4 names the fix\\[2pt]\textcolor{muted}{V8, V10 now}};
  \node[note] at ($(review.south)+(0,-0.15)$) {L0: only a person\\can see it\\[2pt]\textcolor{muted}{V8, V10 at \texttt{8c4a74a}}};
  \node[note] at ($(merge.south)+(0,-0.15)$) {nothing new\\is checked};
  \node[note] at ($(prod.south)+(0,-0.15)$) {found by a user,\\or never};
  % rework loops
  \draw[arr=behaviour] (verify.north) -- ++(0,0.35) -| node[pos=0.25, above]{rework at once; nobody waits} ($(write.north)+(0.45,0)$);
  \draw[arr=missed] (review.north) -- ++(0,0.95) -| node[pos=0.25, above]{rework after a person has waited} ($(write.north)+(-0.45,0)$);
  % direction of the ladder
  \draw[arr=jotblue, line width=1pt] ($(prod.south)+(0,-1.55)$) -- node[below]{each rung climbed moves detection upstream} ($(write.south)+(0,-1.55)$);
\end{tikzpicture}
}
  \caption{The value stream of one change. Each rung moves the point at which a broken rule is
  found further upstream. At commit \texttt{8c4a74a}, V8 and V10 passed the gate and could be seen
  only in review; they are now stopped at the gate. The figure shows position only; we have no
  timing data.}
  \label{fig:valuestream}
\end{figure*}

\paragraph{Root cause analysis} The first seeding run is a small worked example of Lean problem
solving. The problem was concrete: six of 36 seeded mistakes passed the gate. Asking why led to a
cause at the level of the ladder, not of any receiver. The misses were the same two mistakes, V8
and V10, in all three languages. Both passed because every behavioural test passed; every
behavioural test passed because neither mistake changes what the system does; and nothing else
looked, because the structural rules covered which layer may import which, not which layer may
log, notify or decide. The root cause was that behaviour cannot see placement and no rule asserted
placement for side effects or decisions. The countermeasure was two \pattern{Executable Rule}s,
each with a \pattern{Teaching Failure} naming where the code belongs, which is to say two rules
moved from L0 and L2 to L4. The check was a second run, which caught all 42 seeds, including the new requirement
seeds. The standard then
absorbed the countermeasure, in the plan--do--check--act manner of an A3 report
\citep{sobek2008a3}. The V9 fixture episode (Section~\ref{sec:usecases}) followed the same cycle on
a smaller scale, and is the kaizen of a check rather than of the code. What Lean adds to the usual
advice to add a test after a bug is the insistence on the cause: because the same gap appeared in
three languages, the countermeasure was a kind of rule every receiver needed, not a patch to one of
them.

Figure~\ref{fig:a3} writes the episode up as an A3, in the form the Tutors release harness
writes for a release (Section~\ref{sec:tutors}). We include it because it is the ladder's
hallmark on one page. The current condition is where a mistake was found along the stream, the
root cause is the rung it was sitting on, the countermeasure is the rung it climbed to, and the
check is the next run. The form carries no verdict of its own, which is the point. In Tutors the
harness opens an A3 on every failed candidate and its findings inform a confidence score that a
person reads before shipping. In the receivers the countermeasure went straight into the gate.
A team can use the same sheet for confidence or for a hard check, and the ladder is the scale
on which it says which.

\begin{figure*}[t]
  \centering
  \resizebox{0.96\linewidth}{!}{% The first seeding run written up as an A3 report, in the form the Tutors harness writes for a release.
\begin{tikzpicture}[
    font=\scriptsize,
    panel/.style={draw=jotblue!60, line width=0.6pt, fill=white, text width=7.1cm, align=left,
                  inner sep=5pt, anchor=north west},
    head/.style={font=\scriptsize\bfseries, text=jotblue},
    step/.style={draw=#1, line width=0.7pt, fill=#1!7, rounded corners=2pt, minimum height=0.55cm,
                 minimum width=1.5cm, align=center, inner sep=2pt},
    arr/.style={-{Stealth[length=2mm]}, line width=0.7pt, muted},
  ]
  % sheet
  \draw[jotblue, line width=1pt, fill=wash] (0, 0) rectangle (15.4, -10.6);
  \node[anchor=north west, font=\footnotesize\bfseries, text=jotblue] at (0.25, -0.2)
    {A3: six seeded mistakes passed the gate};
  \node[anchor=north east, font=\scriptsize, text=muted] at (15.15, -0.25)
    {receivers at \texttt{8c4a74a}, 21 September 2026; countermeasures measured 29 September};

  % left column
  \node[panel] (bg) at (0.25, -0.8) {
    \textbf{\textcolor{jotblue}{Background}}\\
    The three receivers enforce structure with architecture rules and a short manifest.
    Seeding twelve canonical mistakes into each and running the single verify command,
    the gate caught 30 of 36.};
  \node[panel] (cur) at (0.25, -2.55) {
    \textbf{\textcolor{jotblue}{Current condition}}\\[3pt]
    \begin{tikzpicture}[baseline]
      \node[step=structure] (w) {write};
      \node[step=behaviour, right=0.5cm of w] (g) {gate};
      \node[step=missed, right=0.5cm of g] (r) {review};
      \draw[arr] (w) -- (g);
      \draw[arr] (g) -- node[above, text=missed, font=\tiny]{V8, V10 pass} (r);
    \end{tikzpicture}\\[2pt]
    V8 (a side effect inlined in the service) and V10 (a business rule placed in the route)
    pass every test in Java, Python and TypeScript and reach a human reviewer. The other
    thirty are stopped at the gate.};
  \node[panel] (goal) at (0.25, -5.0) {
    \textbf{\textcolor{jotblue}{Goal}}\\
    Every seeded mistake stopped at the gate, by the earliest mechanism its language allows,
    with no change to application code.};
  \node[panel] (rca) at (0.25, -6.3) {
    \textbf{\textcolor{jotblue}{Root cause (five whys)}}\\
    1. Why did V8 and V10 pass? Every behavioural test passed.\\
    2. Why? Neither mistake changes what the system does.\\
    3. Why did no rule object? The rules say which layer may import which, not which layer
       may log, notify or decide.\\
    4. Why? No rule asserted placement for side effects or decisions.\\
    5. Why in all three languages? Behaviour cannot see placement, and the pattern had no
       structural rule for it. The gap was in the ladder, not in a receiver.};

  % right column
  \node[panel] (cm) at (7.95, -0.8) {
    \textbf{\textcolor{jotblue}{Countermeasures}}\\
    Two \textsc{Executable Rule}s in each receiver: only the events layer may log or notify;
    routes decide nothing. Each carries a \textsc{Teaching Failure} that names where the
    code belongs. V8 and V10 move from L0 and L2 to L4.\\[2pt]
    While there: a requirement \textsc{Closure Check}, with two new seeds (V13, V14) to
    prove it can fail.};
  \node[panel] (chk) at (7.95, -3.45) {
    \textbf{\textcolor{jotblue}{Check}}\\
    Second run: 42 of 42 caught. V7 now found by an architecture rule rather than a
    behavioural test in Python and TypeScript. Cost: 31 to 89 lines of test code and
    80 to 83 words of manifest per receiver. Application code unchanged.};
  \node[panel] (fu) at (7.95, -5.55) {
    \textbf{\textcolor{jotblue}{Follow-up}}\\
    Repeat the pilot against the current receivers with the request in EARS. Raise the
    harness-erosion use cases from L2 to L3. Adopt the release harness on a second product.};
  \node[panel, fill=jotblue!6] (how) at (7.95, -7.3) {
    \textbf{\textcolor{jotblue}{How this sheet is used}}\\
    In Tutors the harness writes an A3 for every release and its findings inform a confidence
    score. In the receivers the countermeasure became part of the gate. The form is the same;
    a team chooses whether it advises or stops the line.};
\end{tikzpicture}
}
  \caption{The first seeding run as an A3 report. One sheet carries the problem, where it was
  found, why, what rung the rule climbed to and what the next run showed. Tutors writes the same
  sheet for a release and reads it as confidence; the receivers made its countermeasure a gate.}
  \label{fig:a3}
\end{figure*}

\paragraph{Supporting mechanisms} Three further Lean mechanisms map onto the ladder more briefly.
The single gate is \emph{jidoka}: the build stops itself on an abnormality, so that no person has
to watch every line, and the manifest's instruction to change the code rather than the test is
the rule that a stop is answered by fixing its cause. No pilot run modified a rule. Jidoka's
older name, automation with a human touch, fits a team of people and agents well: the machine
detects and stops, the person judges. The \pattern{Teaching Failure} is the \emph{andon} signal,
which shows where the problem is and why; the \pattern{Human Boundary} is the andon cord, the
right and the duty to stop and ask a person before changing a dependency, a rule or a schema. The
\pattern{Unrepresentable Mistake} is \emph{poka-yoke} of the preventive kind.

\paragraph{Where the analogy stops} The mapping is close but not exact, and the places where it
breaks tell us something. First, an agent does not learn across shifts. A factory worker is taught the
standard, practises it and in time no longer needs the sheet; an agent is taught afresh every
session. The cost of showing and telling is therefore paid every time, which is the strongest
reason to move a rule from prose into a check, where following it costs no reading. For the same
reason, improvement cannot happen inside the agent. Whatever an agent notices in a session is lost
unless it is written into the repository, so reflection and kaizen run through people and the
files they own. Second, Lean's respect for people becomes the \pattern{Human Boundary}. In Lean the
people who do the work own the standard. Here much of the work is done by agents, so people keep
what needs judgement: what counts as value, how high a rule should climb, and when a check is
wrong. The risk to avoid is the one Lean warned against, reducing people to inspectors of someone
else's output; the ladder gives them the standard and its checks to own instead of every line to
read. Third, the fixture is made of the same material as the work. A worker cannot redesign the jig
in the middle of a shift, but an agent can edit a test, which is why changes to the rules sit
behind the \pattern{Human Boundary}, and why use case UC8 remains at L2 until repository settings
enforce it. Fourth, the ladder can itself become waste. A check is cheap to run but not free to
maintain, and a rule encoded in types far more elaborate than its risk warrants is
over-processing. Reinertsen's caution applies \citep{reinertsen2009principles}, and so do the
findings of Staats and colleagues on uncertain knowledge work \citep{staats2011lean}: standardise
what should repeat and leave the rest free. The ladder applies Lean to a team's conventions, not
to its design.

\subsection{Domain sanctity against ease of use}

A repository can be made arbitrarily hard to damage by making it arbitrarily hard to change.
The ladder resists this in two ways. Each rule climbs only as far as its risk warrants, so the
cost of a rung is paid only where it buys something. Lean would call the alternative
over-processing. And the ladder is compressed from below as
it grows from above: when the receivers gained a requirement file, a scenario convention and
four new rules, their manifests grew by one paragraph each and stayed under 500 words, so the
written standard stayed short enough to be read in every session. The
evidence that ease of use survived is modest but concrete: in the pilot, agents given the full
rule set finished in 151 seconds on average against 121 for the exemplar alone, and used about
6\% more tokens, for a condition in which none departed from the conventions.

\subsection{Where each kind of intent stops}

Each kind of intent sees some mistakes and is blind to others, which is why the three are
complementary rather than redundant. Behaviour cannot see placement: V8 and V10 pass every
scenario because what the system does is unchanged. Structure cannot see meaning: an
architecture rule will not notice that a task may now skip a step. Requirements cannot execute:
without the closure check that ties them to scenarios, they decay into prose, the failure the
manifest's own design avoids. The ladder's contribution is to put all three on the same scale,
so that a team can ask of any rule, in any of the three, how strongly it is asserted.

\subsection{Threats to validity}

\paragraph{Construct} Seeded mistakes are canonical, not observed; real agents may make
mistakes the fourteen do not represent. The TypeScript architecture rules match source text, not
syntax, and could be evaded by code written to avoid them.

\paragraph{Internal} The authors built the receivers, the seeds and the rules together with an
AI coding agent, in the collaborative mode the paper studies. Rules were added after the misses
they close were known; the second seeding run therefore shows that the gap can be closed, not
that the rules generalise to unseen mistakes.

\paragraph{External} One small domain, three mainstream stacks and two models from one vendor
limit generality. The pilot ran one agent per cell and has not yet been repeated against the
current receivers.

\paragraph{Conclusion} Seeding results are deterministic and reproducible from the repository.
Pilot results are suggestive only.
